% !TeX root = ../MAIN.tex
% !TeX encoding = UTF-8

\chapter{Weitere Beispiele}

Hier werden einige weitere Beispiele für importierte Pakete gezeigt.

\section{Todos}

Mit \verb|\todo| kann man Todos erzeugen, wie zum Beispiel hier. \todo{Das ist ein Todo.}

Es ist auch möglich auf noch fehlende Grafiken hinzuweisen, indem man \verb|\missingfigure| angibt.

\missingfigure{Beschreibung der fehlenden Abbildung}

Weiters kann man sich mit \verb|\listoftodos| eine Liste aller Todos ausgeben lassen.

\section{URL}

URLs können mit \verb|\href| oder \verb|\url| eingefügt werden. Ein Beispiel für ersteres ist \href{http://dke.jku.at}{DKE}; für zweiteres \url{https://www.jku.at}.

\section{Listing}

Codeausschnitte können mit Listings eingefügt werden (\verb|\lstinputlisting|).

\lstinputlisting[language=java,label={lst:java},caption={Hello World}]{listings/HelloWorld.java}